\infographicchapter{Ultrasound}{ultrasound.png}
\label{ch:infographic-ultrasound}

\section{Overview}

Diagnostic ultrasound---also called sonography---is one of the safest, most familiar, and most versatile medical imaging tests in healthcare. Instead of X-rays, it uses gentle, high-frequency sound waves (between 2 and 18 MHz, far above human hearing) to create live, moving pictures of organs, soft tissues, and blood flow inside your body. A small handheld probe glides over warm or cool water-based gel on your skin, sending sound pulses inward and listening to the echoes that bounce back from different tissues. A computer turns those echoes into real-time black-and-white images, while special Doppler modes measure blood flow speed and direction.

Because ultrasound works across almost every part of the body, doctors choose the probe frequency to match the depth they need. Lower-frequency probes (2--5 MHz) reach deeper for pregnancy checkups, pelvic exams, liver, gallbladder, kidney, and abdominal scans, and heart ultrasounds (echocardiograms). Higher-frequency probes (7--18 MHz) create crisp close-up views of areas near the skin---including the thyroid, breast, scrotum, neck and leg blood vessels, and joints or tendons---and also guide doctors safely during needle biopsies or fluid drainages.

Ultrasound's biggest advantages are zero X-ray radiation, painless noninvasive scanning, live motion viewing, bedside portability, affordability, and an exceptional ability to tell harmless fluid-filled cysts apart from solid lumps. Its main limits are that picture quality depends strongly on the examiner's skill, it looks at one smaller area at a time rather than the whole body, thick tissue layers can soften deep views, and sound waves cannot travel well through air (such as lungs or bowel gas) or dense bone. Diagnostic ultrasound has an outstanding safety record, and clinicians always keep sound energy as low as reasonably achievable (ALARA), especially during pregnancy and Doppler scans.

\section{Underlying Physical Principles}

\begin{itemize}
    \item \textbf{Sound-Making Crystals in the Probe:} Inside the handheld probe sit tiny piezoelectric crystals. When a brief electrical pulse hits them, they vibrate to send out high-frequency sound waves (2--18 MHz). When returning echoes bounce back from your body, those same crystals vibrate again and turn the echoes back into electrical signals.
    \item \textbf{How Echoes Create Shades of Gray:} Different tissues bounce sound differently based on their density and firmness. Smooth fluid (like urine, bile, or a simple cyst) lets sound pass right through and looks black, with bright tissue behind it. Soft organs look gray and speckled, while air and hard bone bounce nearly all sound back, casting a dark shadow behind them.
    \item \textbf{Timing Depth and Focusing the Beam:} The probe's electronics focus and steer the sound beam smoothly across the target area. By timing how long each echo takes to return, the computer knows exactly how deep each structure sits, and automatically boosts faint echoes from deeper organs so the picture looks evenly bright.
    \item \textbf{Measuring Blood Flow With the Doppler Effect:} Just as an ambulance siren changes pitch as it drives past, sound waves bouncing off moving red blood cells shift slightly in frequency. Color Doppler paints blood flow direction onto the screen, while spectral Doppler graphs exact blood speed to spot narrowed arteries or clots.
    \item \textbf{Advanced Ultrasound Features:} Modern scanners also offer tissue harmonic imaging for cleaner pictures, elastography to measure tissue stiffness (such as liver scarring), and contrast-enhanced ultrasound (CEUS) using tiny harmless gas microbubbles to highlight fine blood flow.
\end{itemize}

\section{Imaging Mechanism}

\begin{enumerate}
    \item \textbf{Choosing the Probe and Applying Gel:} The sonographer picks a probe shaped and tuned for your exam---curved for the abdomen, flat and high-frequency for shallow areas like the thyroid or breast, or specialized for the heart or pelvis---and applies water-based gel (or a sterile cover and sterile gel for internal or needle-guided exams) to remove air gaps.
    \item \textbf{Sweeping the Sound Beam:} The scanner fires rapid, precisely timed sound pulses across the target area while the sonographer glides and tilts the probe, switching smoothly between black-and-white anatomy views, color blood-flow Doppler, or stiffness modes as needed.
    \item \textbf{Listening to Returning Echoes:} Returning echoes vibrate the probe's crystals into tiny electrical signals, which are boosted, balanced for depth, and converted into digital data.
    \item \textbf{Building Live Pictures and Flow Graphs:} The computer cleans up speckle noise, sharpens tissue edges, and calculates live blood-flow colors, velocity graphs, or stiffness maps in real time.
    \item \textbf{Viewing, Measuring, and Saving:} Live moving pictures appear immediately on the screen alongside safety gauges. The sonographer measures key structures and saves still images and short video clips for the radiologist to interpret.
\end{enumerate}

\section{Key Technical Parameters}

\noindent
\begin{tabularx}{\textwidth}{@{}l X@{}}
\toprule
\textbf{Parameter} & \textbf{Typical Value / Description} \\
\midrule
Physical Principle & Pulse-echo reflection of high-frequency sound waves at tissue boundaries, plus Doppler blood-flow and stiffness modes \\
Energy / Signal Source & Handheld piezoelectric ultrasound probe (curved, flat linear, heart phased-array, or internal) \\
Typical Frequency / Energy & 2--5 MHz (abdomen, pelvis, pregnancy, heart) and 7--18 MHz (thyroid, breast, scrotum, blood vessels, joints/tendons) \\
Spatial Resolution & About 0.2--1.5 mm (finer detail with high-frequency shallow probes) \\
Typical Scan Duration & 15--45 minutes \\
Contrast Agent & None routinely; optional intravenous microbubble contrast for specialized CEUS exams \\
Ionizing Radiation & None (0 mSv) \\
\bottomrule
\end{tabularx}

\section{Clinical Usages}

\subsection{Primary Usages}

\begin{itemize}
    \item \textbf{Pregnancy and Women's Health:} The go-to first exam for confirming pregnancy, checking due dates, monitoring baby growth and anatomy, checking the placenta and fluid, and evaluating uterine fibroids, ovarian cysts, or pelvic pain.
    \item \textbf{Abdomen, Liver, and Kidneys:} The first-choice test for gallstones, gallbladder inflammation, blocked bile ducts, fatty liver or liver spots, swollen kidneys, kidney masses, and screening for an enlarged abdominal aorta (aneurysm).
    \item \textbf{Blood Vessels and Clots:} Checking leg or arm veins for blood clots (deep vein thrombosis), screening neck (carotid) arteries for narrowing, checking leg circulation, and verifying blood flow in transplanted organs.
    \item \textbf{Shallow Organs and Children:} Evaluating thyroid nodules, breast lumps (telling fluid cysts apart from solid lumps, especially in younger or dense breasts), sudden scrotal pain (testicular torsion or infection), and common childhood conditions like pyloric stenosis, intussusception, or infant hip development.
    \item \textbf{Heart Ultrasound (Echocardiography):} Checking heart chamber size, pumping strength, leaky or stiff heart valves, fluid around the heart, and congenital heart conditions.
\end{itemize}

\subsection{Secondary Usages}

\begin{itemize}
    \item \textbf{Live Needle Guidance:} Guiding needle biopsies (thyroid, breast, liver, lymph nodes), fluid drainages from the chest or abdomen, abscess drains, joint injections, and tricky IV placements in real time.
    \item \textbf{Joints, Tendons, and Muscles:} Checking rotator cuff tears, tendon inflammation, bursitis, ligament sprains, pinched nerves (like carpal tunnel syndrome), and surface lumps while you move the joint.
    \item \textbf{Tissue Stiffness Checks (Elastography):} Measuring liver stiffness to check for scarring without a biopsy, and adding stiffness clues for breast or thyroid nodules.
    \item \textbf{Contrast-Enhanced Ultrasound (CEUS):} Using tiny, kidney-safe gas microbubbles to clarify liver or kidney spots and check aortic stent grafts for leaks without X-ray radiation.
\end{itemize}

\section{Advantages}

\begin{itemize}
    \item \textbf{Zero Radiation and Proven Safety:} Uses gentle sound waves rather than X-rays, making it the safest first-line scan during pregnancy, for children, and for repeat follow-ups.
    \item \textbf{Live Motion and Blood Flow:} Shows organs moving in real time---like a beating heart, a moving baby, or gliding tendons---measures blood flow with Doppler, and guides needles safely.
    \item \textbf{Clearly Separates Cysts From Solid Lumps:} Reliably tells harmless fluid-filled cysts apart from solid or complex masses in the breast, thyroid, ovaries, kidneys, and liver.
    \item \textbf{Widely Available, Portable, and Affordable:} Costs far less than CT or MRI, is available almost everywhere, and portable machines can roll right to a hospital bedside or emergency room.
    \item \textbf{Noninvasive and Usually Dye-Free:} Most scans need no IV line or contrast dye; when microbubble contrast is used, it is breathed out through the lungs and is safe for the kidneys.
\end{itemize}

\section{Disadvantages}

\begin{itemize}
    \item \textbf{Depends Strongly on Examiner Skill:} Picture quality and accuracy rely heavily on the sonographer's technique in positioning the probe and tuning the settings, so results can vary more between operators than automated CT or MRI scans.
    \item \textbf{Blocked by Air, Bone, and Deep Tissue Layers:} Sound bounces off air and bone, making it hard to image air-filled lungs, the adult brain, organs hidden behind bowel gas, or areas behind thick bone. Heavy body weight can also soften deep images.
    \item \textbf{Smaller Viewing Window and Echo Tricks:} Shows one smaller region at a time rather than a whole-body overview, and harmless echo illusions (like shadowing or mirror images) can occasionally mimic real findings.
    \item \textbf{Keeping Sound Energy Gentle:} High sound power---especially prolonged pulsed Doppler in early pregnancy---can cause slight local warming, so clinicians strictly follow ALARA safety limits.
    \item \textbf{Not Right for Every Question:} Cannot replace CT or MRI for full-body cancer staging, sudden brain bleeding, complex bone fractures, or deep lung disease.
\end{itemize}

\section{Contraindications}

\subsection{Absolute}

\begin{itemize}
    \item \textbf{None for Routine Skin-Surface Ultrasound:} Standard grayscale and Doppler ultrasound over intact skin has no strict medical contraindications.
    \item \textbf{Allergy to Microbubble Contrast (CEUS Only):} Intravenous microbubble contrast must not be used if you have a known allergy to the microbubble agent.
\end{itemize}

\subsection{Relative}

\begin{itemize}
    \item \textbf{Internal Scans and Open Wounds:} Internal (transvaginal or transrectal) scans are avoided if there is a severe local blockage, very recent pelvic or rectal surgery, or if a patient declines. Scanning over open wounds or fresh incisions uses a sterile probe cover and sterile gel.
    \item \textbf{Biopsies and Contrast Precautions:} Uncorrected severe bleeding disorders are a caution before ultrasound-guided needle biopsies or deep drainages, and severe unstable heart or lung conditions warrant extra care before microbubble contrast.
\end{itemize}

\section{Patient Safety Guidelines}

\begin{itemize}
    \item \textbf{ALARA Safety Limits:} Following AIUM and WFUMB guidelines, sonographers use the lowest sound power and shortest scan time needed, watching on-screen Thermal Index (TI) and Mechanical Index (MI) gauges---especially in early pregnancy, pulsed Doppler, and eye scans (TI at or below 1.0, MI at or below 0.23).
    \item \textbf{Probe Cleaning and Infection Control:} Probes and cords are wiped with disinfectant between every routine skin scan. Internal probes, scans over broken skin, and needle-guided procedures require high-level disinfection plus a fresh sterile cover and single-use sterile gel.
    \item \textbf{Microbubble Contrast Safety (CEUS):} Before CEUS, staff check for allergies and heart-lung stability, use low-energy contrast settings, and keep emergency care supplies nearby.
    \item \textbf{Equipment Checks:} Staff inspect probe faces and cables before each scan for cracks or overheating to keep you safe from electrical or heat irritation.
    \item \textbf{Comfort in Late Pregnancy and Needle Safety:} In mid-to-late pregnancy, lying slightly tilted on your left side avoids pressure on major abdominal blood vessels. During biopsies, the clinician watches the needle tip continuously on screen.
\end{itemize}

\section{Prerequisite Tests}

\begin{itemize}
    \item \textbf{No Routine Blood Tests for Standard Scans:} Regular grayscale and Doppler ultrasounds are noninvasive and dye-free, so no blood work is needed beforehand.
    \item \textbf{Simple Exam-Specific Prep:} Check your appointment instructions---abdominal and gallbladder scans usually require 6--8 hours of fasting, while pelvic and early pregnancy scans often ask for a comfortably full bladder.
    \item \textbf{Pregnancy Test (When Helpful):} For pelvic pain or unusual bleeding in patients who could be pregnant, a quick urine or blood pregnancy test helps the team tailor the scan and interpret early pregnancy findings.
    \item \textbf{Blood-Clotting Checks (For Biopsies Only):} Before an ultrasound-guided needle biopsy or deep fluid drainage, your doctor reviews blood clotting tests (INR and platelets) and any blood-thinning medicines.
    \item \textbf{Allergy and Heart-Lung Check (For CEUS Only):} Before microbubble contrast, staff check for past reactions and heart-lung stability; no kidney blood test is needed because the gas is breathed out through the lungs.
\end{itemize}

\section{Preparing for the Scan}

\begin{itemize}
    \item \textbf{Fasting for Abdomen and Gallbladder Scans:} For liver, gallbladder, pancreas, or abdominal aorta scans, avoid food, fizzy drinks, and chewing gum for 6--8 hours beforehand so the gallbladder stays full and bowel gas stays low. Small sips of plain water with essential pills are usually fine.
    \item \textbf{Full Bladder for Pelvic and Early Pregnancy Scans:} For lower-belly pelvic, bladder, or early pregnancy scans, drink about 24--32 ounces (700--950 mL) of water finishing 45--60 minutes before your visit and do not empty your bladder, as a full bladder acts like a clear window. (If an internal transvaginal scan is planned, you will empty your bladder right before that part.)
    \item \textbf{Little to No Prep for Shallow or Blood-Vessel Scans:} Thyroid, breast, scrotal, joint, neck artery, and leg vein scans need no fasting or full bladder---just skip thick lotions, oils, or powders on the area being scanned.
    \item \textbf{Easy-Access Clothing:} Wear loose, comfortable two-piece clothing so the area can be uncovered easily, and remove necklaces or jewelry near the scan site.
    \item \textbf{Biopsy or Contrast Prep:} If you are scheduled for a needle biopsy, drainage, or microbubble contrast scan, follow your clinic's specific instructions on blood thinners, fasting if sedation is planned, IV placement, and having a ride home.
\end{itemize}

\section{During Scan}

\begin{itemize}
    \item \textbf{Comfortable Positioning:} You lie on a padded exam bed---usually on your back, turning slightly onto your side, or sitting upright depending on the body part (with a gentle left-side tilt in late pregnancy)---and stay relaxed while the sonographer glides the probe.
    \item \textbf{Simple Breathing Cues:} You breathe normally for most of the exam. For liver, gallbladder, or kidney views, the sonographer may ask you to take a deep breath and hold it for 5--10 seconds so the organs drop below your ribs for a clearer view.
    \item \textbf{What You Will Feel and Hear:} You will feel cool or warmed gel on your skin and gentle-to-firm gliding pressure from the smooth probe (which can feel briefly snug over a full bladder or sore spot). During Doppler checks, you may hear a gentle rhythmic whooshing sound of your blood flow.
    \item \textbf{Speak Up Anytime:} The sonographer stays right beside you the entire time. Tell them immediately if probe pressure hurts, if lying flat makes you dizzy or short of breath, or if your full bladder feels painfully tight.
    \item \textbf{Time and Pacing:} Most diagnostic ultrasounds take 15--45 minutes depending on how many areas and blood-flow checks are needed, with zero recovery time afterward for standard scans.
\end{itemize}

\section{Normal Appearance}

\begin{itemize}
    \item \textbf{Smooth, Even Solid Organs:} A healthy liver, spleen, thyroid, or testicle looks smooth, finely speckled, and even in shade with neat outer borders and no lumps or hard shadows.
    \item \textbf{Clean Fluid Spaces:} A fasting gallbladder appears as a thin-walled, pear-shaped black sac without bright stones or surrounding fluid, and the bladder fills cleanly with clear black urine.
    \item \textbf{Healthy Kidneys and Blood Vessels:} Normal kidneys show smooth outer tissue and a bright center without backed-up fluid. Healthy veins flatten easily under gentle probe pressure, and arteries show smooth, steady blood-flow waves.
    \item \textbf{Healthy Tendons and Pregnancy:} Healthy tendons show neat, bright parallel fibers, while a normal early pregnancy shows an appropriately sized sac, expected baby growth, healthy fluid levels, and a regular heartbeat.
\end{itemize}

\section{Observing for Abnormal Patterns}

\begin{itemize}
    \item \textbf{Simple Fluid-Filled Cysts:} A smooth, round, completely black spot with a bright glow directly behind it is the classic sign of a harmless fluid-filled cyst in the kidney, liver, breast, thyroid, or ovary.
    \item \textbf{Bright Stones and Gallbladder Inflammation:} Bright white spots casting a sharp dark shadow behind them identify gallstones, kidney stones, or hard calcium deposits. A thickened gallbladder wall (over 3 mm), fluid around it, and tenderness under the probe point to acute gallbladder inflammation (cholecystitis).
    \item \textbf{Suspicious Solid Lumps and Fatty Liver:} Jagged, taller-than-wide dark solid nodules with tiny bright calcium specks or tangled internal blood flow raise concern for cancer in the breast or thyroid. An unusually bright liver that dims deeper views points to fatty liver (hepatic steatosis).
    \item \textbf{Blood-Flow Problems on Doppler:} A vein that will not flatten under gentle pressure and lacks blood flow confirms a blood clot (DVT). Very fast, turbulent blood flow pinpoints a narrowed artery (such as in the neck), while missing blood flow in a painful testicle or ovary flags urgent twisting (torsion).
    \item \textbf{Common Echo Illusions:} Bowel gas or surgical clips can create bright repeating streaks, reflections off the diaphragm can mirror a liver spot into the chest, and faint echo blur can make a clear cyst look like it has debris inside unless checked from another angle.
\end{itemize}

\section{Dealing with Uncertainties}

\begin{itemize}
    \item \textbf{Gas, Bone, and Deep-Tissue Barriers:} Heavy body weight softens deep sound waves, and overlying bowel gas, lungs, or ribs can block views of the pancreas, deep pelvis, or abdominal aorta.
    \item \textbf{Technique and Echo Illusions:} Image clarity depends strongly on probe angle, frequency choice, and settings; echo artifacts can sometimes mimic debris inside cysts or exaggerate artery narrowing.
    \item \textbf{Prep Challenges:} An empty bladder during a pelvic scan or a shrunk gallbladder after a recent meal can make the view unreadable or falsely look like a thickened gallbladder wall.
    \item \textbf{Borderline Lumps:} Some complex cysts, liver or kidney spots, or thyroid and ovarian nodules share features of both harmless and serious conditions, so standard ultrasound alone cannot always give a final answer.
    \item \textbf{How Doctors Clarify Results:} Clinics may repeat a scan after proper fasting or drinking water, shift your position or try a different probe angle, add Doppler, stiffness checks (elastography), or microbubble contrast (CEUS), or recommend CT, MRI, a follow-up check, or a quick needle biopsy.
\end{itemize}

\section{Augmenting Methods}

\begin{itemize}
    \item \textbf{Contrast-Enhanced Ultrasound (CEUS) and Elastography:} Microbubble contrast clarifies liver and kidney spots in real time without harming the kidneys, while elastography measures liver scarring or lump stiffness.
    \item \textbf{CT and MRI Scans:} Overcome bowel gas, bone, and deep-tissue limits to give wide-angle 3D views for cancer staging, pancreatic disease, and detailed soft-tissue checks.
    \item \textbf{Ultrasound-Guided Needle Biopsy:} Uses live ultrasound guidance to take a tiny, accurate tissue sample from uncertain thyroid, breast, liver, or lymph node spots with minimal discomfort.
    \item \textbf{Targeted Blood Tests:} Pairing ultrasound with pregnancy hormones (beta-hCG), liver tests, blood-clot checks (D-dimer), PSA, or CA-125 sharpens the diagnosis.
    \item \textbf{Standardized Risk Scores and AI Support:} Clear scoring systems (BI-RADS for breast, TI-RADS for thyroid, LI-RADS for liver, O-RADS for ovaries) and modern AI measurement tools keep reporting consistent and evidence-based.
\end{itemize}

\section{Comparison with Alternative Modalities}

\begin{itemize}
    \item \textbf{Versus CT Scan:} CT gives wide, rapid 3D views unaffected by bowel gas or body size (better for major trauma and cancer staging), while ultrasound uses zero radiation and no kidney-stressing dye and shines for pregnancy, gallstones, blood clots, and shallow organs.
    \item \textbf{Versus MRI:} MRI gives unmatched deep soft-tissue detail (for brain, spine, or complex liver and pelvic masses), while ultrasound is much faster, far less expensive, widely available, and ideal for live motion and needle guidance.
    \item \textbf{Versus Bedside Quick Scans (POCUS):} Focused bedside ultrasound answers urgent, single questions in seconds during emergencies, whereas a full diagnostic ultrasound provides thorough multi-organ views, detailed Doppler measurements, and formal radiology reporting.
\end{itemize}

\section{Follow-up Steps}

\begin{itemize}
    \item \textbf{Right After the Exam:} Standard diagnostic ultrasound needs zero recovery time---wipe off the gel and go right back to your day. After microbubble contrast (CEUS), you may rest briefly for observation; after a biopsy or drainage, staff monitor your vital signs and puncture site briefly before you head home.
    \item \textbf{Clear Reporting and Urgent Alerts:} Radiologists use standardized scoring systems (BI-RADS, TI-RADS, O-RADS, LI-RADS) and call your doctor immediately if an urgent finding---like a blood clot, ectopic pregnancy, torsion, or leaking aneurysm---needs fast care.
    \item \textbf{Next-Step Imaging or Biopsy:} If a scan is limited by gas or shows a borderline spot, your doctor may order a targeted CT, MRI, or a simple ultrasound-guided needle biopsy for a definitive answer.
    \item \textbf{Routine Checkups Over Time:} For harmless or low-risk findings, your doctor may simply schedule a follow-up ultrasound---such as tracking a small abdominal aneurysm, rechecking a thyroid nodule or ovarian cyst, or monitoring baby growth in later pregnancy.
\end{itemize}

\section{What You Should Know About Ultrasound?}

\begin{itemize}
    \item \textbf{Painless and Noninvasive:} A routine ultrasound involves only cool or warmed water-based gel and gentle-to-firm gliding of a smooth handheld probe across your skin. It is painless, though pressing over a full bladder or tender area can feel briefly snug.
    \item \textbf{Completely Radiation-Free:} Diagnostic ultrasound uses safe sound waves instead of X-rays and normally needs no contrast dye, making it the safest first choice during pregnancy, for children, and for repeat checkups.
    \item \textbf{Quick Visit and Clear Turnaround:} Most scans take 15--45 minutes with no recovery time afterward. Urgent scans are read quickly, while routine results typically reach your doctor within 1--3 business days.
    \item \textbf{Why Your Doctor Orders It:} Ultrasound shows live movement of organs, babies, and blood flow in real time and clearly separates harmless fluid-filled cysts from solid lumps---helping diagnose gallstones, kidney stones, blood clots, thyroid or breast nodules, and pregnancy milestones.
    \item \textbf{Most Important Thing to Tell Your Care Team:} Confirm ahead of time whether your scan requires fasting (for abdomen or gallbladder scans) or a full bladder (for pelvic or early pregnancy scans), and speak up during the exam if probe pressure hurts or lying flat feels uncomfortable.
\end{itemize}

\section{Key References}

\begin{itemize}
    \item American Institute of Ultrasound in Medicine (AIUM) and American College of Radiology (ACR). \textit{AIUM--ACR--SPR--SRU Practice Parameters for Diagnostic Ultrasound Examinations}.
    \item World Federation for Ultrasound in Medicine and Biology (WFUMB). \textit{Safety Guidelines for Diagnostic Ultrasound, Elastography, and Contrast-Enhanced Ultrasound}.
    \item Rumack CM, Levine D. \textit{Diagnostic Ultrasound}, 5th ed. Philadelphia: Elsevier.
    \item Carovac A, Smajlovic F, Junuzovic D. Application of ultrasound in medicine. \textit{Acta Informatica Medica}. 2011;19(3):168--171.
\end{itemize}
